% P16-C4: complete finite soundness of the three do-calculus rules.
\section{Soundness of the three do-calculus rules}
\label{sec:p16-do-calculus-soundness}

Pearl's three rules transform probabilities involving interventions and
observations under explicit graphical conditions
\cite[Theorem~3, pp.~676--678; Appendix, pp.~686--687]{Pearl1995CausalDiagrams}.
This section proves their soundness for the finite structural models
already defined. The rules apply to arbitrary disjoint sets of vertices,
and their conditional expressions are used only on positive events.
Latent endogenous variables remain present in the full DAG.

\begin{hbdefinition}{def:p16-do-calculus-data}{Index sets, graph modifications, and conditional laws}
Fix a finite structural model \(\mathcal M\) on \(G=(N,E)\)
with mutually independent disturbances. In this section \(X,Y,Z,W\)
are pairwise disjoint subsets of \(N\), rather than names of single
random variables. Their assignment tuples are respectively
\(x\in S_X,y\in S_Y,z\in S_Z,w\in S_W\).
Empty index sets have the single empty assignment.

Write \(G_{\overline A}\) for deletion of all arrows into \(A\),
and \(G_{\underline B}\) for deletion of all arrows out of \(B\).
The notation \(G_{\overline A,\underline B}\) performs both deletions.
Two overlined sets delete incoming arrows into their union. These
operations are specified on the original edge set, so their order
does not change the result. Put \(K=G_{\overline X}\).

Let \(p_x\) be the full assignment law after replacing the equations
at \(X\) by \(x\). For an event \(D\subseteq S_N\), put
\(p_x(D)=\sum_{v\in D}p_x(v)\). For example,
\[
 p_x(y,z,w)=\mathbb P(V_Y^{X=x}=y,V_Z^{X=x}=z,V_W^{X=x}=w).
\]
Omitted coordinates are summed out. For positive denominators define
\[
 p_x(y\mid z,w)=\frac{p_x(y,z,w)}{p_x(z,w)},\qquad
 p_x(y\mid w)=\frac{p_x(y,w)}{p_x(w)}.
\]
The notation \(p_{x,z}\) means simultaneous equation replacement
at \(X\cup Z\) by the indexed tuple combining \(x,z\).
Thus an intervened tuple occurs in the subscript; an observed tuple
occurs after the conditional bar. When \(X\) is empty, write
\(p\) for \(p_x\). The observational probability of the original
event \(V_X=x\) need not be positive.

Pearl's source notation writes
\(\operatorname{pr}(y\mid\check x,z,w)\) for \(p_x(y\mid z,w)\),
\(\operatorname{pr}(y\mid\check x,\check z,w)\) for
\(p_{x,z}(y\mid w)\), and
\(\operatorname{pr}(y\mid\check x,w)\) for \(p_x(y\mid w)\).
The check mark there designates an intervened tuple; it does not
designate an additional observed random variable.

For action deletion define exactly
\[
 Z(W)=Z\setminus\operatorname{an}_{K}(W).
\]
These are the proposed intervention vertices that are not ancestors of
any conditioned vertex after the baseline intervention at \(X\).
Since vertices of \(X\) have no incoming arrows in \(K\) and
\(Z\cap X=\varnothing\), this set also equals
\(Z\setminus\operatorname{an}_{K}(X\cup W)\).
In particular, \(Z(\varnothing)=Z\).
\index{do-calculus}\index{action deletion!ancestor exception}
\end{hbdefinition}

\begin{hbdefinition}{def:p16-independent-regimes}{Independent finite regime variables}
For fixed \(x,z\), start with the model after intervention \(X=x\).
Adjoin one fresh binary root \(F_i\) for each \(i\in Z\),
independent fair bits jointly independent of all original disturbances.
The original equation at \(i\in Z\) is used when \(F_i=0\),
and is replaced by the constant value \(z_i\) when \(F_i=1\).
Equations at all other original vertices are retained. The allowed graph
\(\widetilde K\) contains \(K\) and the new edges \(F_i\to i\).
Write \(\mathcal F_Z=\{F_i:i\in Z\}\) for the new vertex set,
\(F=(F_i)_{i\in Z}\) for its random tuple, and \(\widetilde V\)
for the solution on the original vertices. Denote the augmented
outcome probability by \(\widetilde{\mathbb P}\).
The regime variables are auxiliary model data for the proof; they are
not observed treatment choices in the original population.
\end{hbdefinition}

\begin{hblemma}{thm:p16-regime-law}{proved=annex:p16.regime-law}{Conditioning on a regime recovers its intervention law}
The augmented construction is a finite acyclic structural model with
mutually independent disturbances. For a bit tuple \(e\in\{0,1\}^Z\),
let \(T_e=\{i\in Z:e_i=1\}\). For every event \(D\subseteq S_N\),
\[
 \widetilde{\mathbb P}(\widetilde V\in D,F=e)
 =2^{-|Z|}\,p_{x,z_{T_e}}(D).
\]
Consequently the all-zero regime gives \(p_x\), and the all-one
regime gives \(p_{x,z}\). The equality includes events of zero mass.
\end{hblemma}

\begin{hblemma}{thm:p16-regime-graph-bridges}{proved=annex:p16.regime-graph-bridges}{The graphical bridges for exchange and action deletion}
Put \(C=X\cup W\), and use the augmented graph just defined.
The following two implications hold:
\begin{enumerate}
\item If \(Y\mathrel{\perp_{K_{\underline Z}}}Z\mid C\), then
      \(Y\mathrel{\perp_{\widetilde K}}\mathcal F_Z\mid(C\cup Z)\).
\item If \(Y\mathrel{\perp_{K_{\overline{Z(W)}}}}Z\mid C\), then
      \(Y\mathrel{\perp_{\widetilde K}}\mathcal F_Z\mid C\).
\end{enumerate}
The second implication retains incoming arrows at members of
\(Z\cap\operatorname{an}_K(W)\).
\end{hblemma}

\begin{hbtheorem}{thm:p16-do-rule-one}{proved=annex:p16.do-rule-one}{Rule 1: insertion and deletion of observations}
If
\[
 Y\mathrel{\perp_{G_{\overline X}}}Z\mid(X\cup W)
 \quad\text{and}\quad p_x(z,w)>0,
\]
then \(p_x(w)>0\) and, for every response assignment \(y\),
\[
 p_x(y\mid z,w)=p_x(y\mid w).
\]
\end{hbtheorem}

\begin{hbtheorem}{thm:p16-do-rule-two}{proved=annex:p16.do-rule-two}{Rule 2: action and observation exchange}
If
\[
 Y\mathrel{\perp_{G_{\overline X,\underline Z}}}Z\mid(X\cup W)
 \quad\text{and}\quad p_x(z,w)>0,
\]
then \(p_{x,z}(w)>0\) and, for every \(y\),
\[
 p_{x,z}(y\mid w)=p_x(y\mid z,w).
\]
The support assumption on the observational side supplies positivity
on the intervention side by pointwise consistency.
\end{hbtheorem}

\begin{hbtheorem}{thm:p16-do-rule-three}{proved=annex:p16.do-rule-three}{Rule 3: insertion and deletion of actions}
If
\[
 Y\mathrel{\perp_{G_{\overline X,\overline{Z(W)}}}}Z\mid(X\cup W),
 \qquad p_x(w)>0,\qquad p_{x,z}(w)>0,
\]
then, for every \(y\),
\[
 p_{x,z}(y\mid w)=p_x(y\mid w).
\]
Here \(Z(W)\) is computed in \(G_{\overline X}\), before its
additional incoming-arrow deletion. Both support conditions are stated
because the proposed action can change the probability of the
conditioning event.
\end{hbtheorem}

These are equalities under their stated hypotheses, so each permits
the indicated replacement in either direction. A proposed symbolic
derivation must satisfy the graph and support hypotheses of every
step. Soundness does not claim that these rules find every identifiable
effect, or supply an algorithm for finding a suitable sequence.

\subsection{A sequence using all three rules}

\begin{hbdefinition}{def:p16-do-sequence-model}{A noisy chain with an independent observation}
For this example use named binary vertices \(A,B,C,D\), with
graph \(A\to B\to C\) and isolated \(D\). Let \(U_A,U_D\)
be fair bits, let \(\mathbb P(U_B=1)=1/4\), and let
\(\mathbb P(U_C=1)=1/5\), all mutually independent. Set
\[
 A=U_A,\qquad B=A\oplus U_B,\qquad
 C=B\oplus U_C,\qquad D=U_D.
\]
The finite outcome space is the four-bit product with these product
masses. Define \(\alpha(0)=3/4,\alpha(1)=1/4\) and
\(\beta(0)=4/5,\beta(1)=1/5\).
When a subscript uses a named variable, such as \(p_a\), it denotes
intervention at that variable with the displayed value.
\end{hbdefinition}

\begin{hbtheorem}{thm:p16-do-rule-sequence}{proved=annex:p16.do-rule-sequence}{A fully justified four-step rule sequence}
For every \(a,b,c,d\in\{0,1\}\) in the preceding model,
\[
\begin{aligned}
 p_a(c\mid b,d)
 & =p_a(c\mid b) &&\text{by Rule 1},\\
 & =p_{a,b}(c) &&\text{by Rule 2},\\
 & =p_b(c) &&\text{by Rule 3},\\
 & =p(c\mid b) &&\text{by Rule 2}\\
 & =\beta(c\oplus b).
\end{aligned}
\]
All required conditioning events have positive mass. The response-one
probability is \(1/5\) at \(b=0\) and \(4/5\) at \(b=1\),
independent of \(a,d\) after conditioning on \(B=b\).
\end{hbtheorem}

\subsection{Failed graph and support conditions}

\begin{hbdefinition}{def:p16-do-ancestor-countermodel}{A proposed action that affects the conditioning variable}
Name binary vertices \(U,Z,W,Y\), with equations
\[
 U=U_U,\qquad Z=U\oplus U_Z,\qquad
 W=Z\oplus U_W,\qquad Y=U.
\]
The root disturbance \(U_U\) is a fair bit, \(U_Z,U_W\) are
independent bits of probability \(1/4\) for one, and the unused
disturbance for \(Y\) is constant; all are mutually independent.
The graph has edges \(U\to Z\), \(Z\to W\), and \(U\to Y\).
The three-bit outcome \((u,a,b)\) has mass
\(\frac12\alpha(a)\alpha(b)\), with \(\alpha\) from
\cref{def:p16-do-sequence-model}.
\end{hbdefinition}

\begin{hbtheorem}{thm:p16-do-failed-conditions}{proved=annex:p16.do-failed-conditions}{The ancestor exception and support assumptions cannot be dropped}
In \cref{def:p16-do-ancestor-countermodel}, for either \(z=0,1\),
\[
 p(Y=1\mid W=1)=5/8,\qquad p_z(Y=1\mid W=1)=1/2,
\]
and both conditioning events are positive. Since \(Z\) is an
ancestor of \(W\), Rule 3 has \(Z(W)=\varnothing\) and its
graph condition fails. Incorrectly deleting all incoming arrows into
\(Z\) instead would make the graphical test pass and assert this
false equality.

For a separate support example, take binary \(Z,Y\) with graph
\(Z\to Y\), singleton disturbances, and equations \(Z=0,Y=Z\).
The graph condition of Rule 2 holds after deleting the outgoing edge
from \(Z\), but \(p(Z=1)=0\). The intervention gives
\(p_1(Y=1)=1\), while \(p(Y=1\mid Z=1)\) is undefined.
Thus a successful graph test alone does not license a conditional
expression on an absent event.

For Rule 3's intervention-side support, let \(Z,Y\) be independent
fair bits and let \(W=Z\), with graph \(Z\to W\) and isolated
\(Y\). Its correct graph test passes, but \(p(W=1)=1/2\) and
\(p_0(W=1)=0\). Thus \(p(Y=1\mid W=1)=1/2\), whereas
\(p_0(Y=1\mid W=1)\) is undefined.
\end{hbtheorem}

The section proves all three rules in the specified finite model class.
\Cref{sec:p16-finite-counterfactuals} next supplies same-unit counterfactual
reasoning and the binary potential-outcome identification results.
Continuous conditional-kernel versions, latent-projection graph
formalisms, identification completeness, and identification algorithms
require separate statements and proofs.
